\documentclass[10pt,a4paper]{amsart}

\usepackage[margin=0.88in]{geometry}
\usepackage[T1]{fontenc}
\usepackage{lmodern}
\usepackage{microtype}
\usepackage{amsmath,amssymb,amsthm,mathtools}
\usepackage{enumitem}
\usepackage{booktabs}
\usepackage{xcolor}
\usepackage[colorlinks=true,linkcolor=blue!60!black,
  citecolor=blue!60!black,urlcolor=blue!60!black]{hyperref}

\allowdisplaybreaks
\setlength{\jot}{6pt}
\setlength{\abovedisplayskip}{5pt}
\setlength{\belowdisplayskip}{5pt}
\setlength{\abovedisplayshortskip}{3pt}
\setlength{\belowdisplayshortskip}{3pt}
\setlength{\textfloatsep}{8pt plus 2pt minus 2pt}
\setlength{\floatsep}{6pt plus 2pt minus 2pt}
\setlength{\intextsep}{6pt plus 2pt minus 2pt}
\setlist{itemsep=2pt,topsep=3pt,parsep=0pt,partopsep=0pt}

% ─── theorem environments ───
\theoremstyle{plain}
\newtheorem{theorem}{Theorem}[section]
\newtheorem{lemma}[theorem]{Lemma}
\newtheorem{proposition}[theorem]{Proposition}
\newtheorem{corollary}[theorem]{Corollary}
\newtheorem{conjecture}[theorem]{Conjecture}

\theoremstyle{definition}
\newtheorem{definition}[theorem]{Definition}

\theoremstyle{remark}
\newtheorem{remark}[theorem]{Remark}

% ─── custom macros ───
\newcommand{\R}{\mathbb{R}}
\newcommand{\PnR}[1][n]{\mathcal{P}^{\R}_{#1}}
\newcommand{\fp}{\boxplus_n}
\DeclareMathOperator{\tr}{tr}
\DeclareMathOperator{\disc}{disc}
\DeclareMathOperator{\Bez}{Bez}
\DeclareMathOperator{\rank}{rank}


% ───────────────────────────────────────────────────────────
\begin{document}

\title[The Finite Free Stam Inequality]%
{The Finite Free Stam Inequality:\\
Structural Identities, the Gaussian-Input Case,
and Reduction Principles}

\subjclass[2020]{46L54 (primary); 94A17, 26C10, 15A42 (secondary)}
\keywords{finite free convolution, Stam inequality, Fisher information,
real-rooted polynomials, Marcus--Spielman--Srivastava convolution}

\date{\today}

\begin{abstract}
Let $\PnR$ denote the set of monic, degree-$n$, real-rooted
polynomials and let $\fp$ be the Marcus--Spielman--Srivastava
finite free additive convolution.
For $r\in\PnR$ with simple roots
$\lambda_1<\cdots<\lambda_n$, the
\emph{finite free Fisher information} is
$\Phi_n(r):=\sum_{i=1}^n V_i^2$, where
$V_i:=\sum_{j\ne i}(\lambda_i-\lambda_j)^{-1}$.
We prove the \emph{finite free Stam inequality}
\[
  \frac{1}{\Phi_n(p\fp q)}\;\ge\;
  \frac{1}{\Phi_n(p)}+\frac{1}{\Phi_n(q)},
  \qquad p,q\in\PnR,
\]
for all $n\ge 2$, confirming a conjecture of D.~Shlyakhtenko.
The proof, due to Garza~Vargas, Srivastava, and
Stier~\cite{GVS}, proceeds in three steps:
(1)~a Jacobian identity relating score vectors under
convolution, derived from the commutation of $\fp$
with differential operators;
(2)~a contraction inequality for the Jacobian of the
root map, established via variance additivity and
a convexity result of Bauschke et al.~\cite{BGLS}
on hyperbolic polynomials;
(3)~Blachman's classical argument adapted to the
finite free setting.
We also develop structural identities
($\Phi_n=2\mathcal{R}=\tr(L)$, Bezoutian representation),
the Hermite heat equation and root ODE,
and give independent proofs for $n\le 3$ and
for Gaussian inputs at all~$n$.
\end{abstract}

\maketitle

% ═══════════════════════════════════════════════════════════
\section{Introduction}
\label{sec:intro}
% ═══════════════════════════════════════════════════════════

The classical Stam inequality~\cite{stam} states that
for independent continuous random variables $X,Y$ with
finite Fisher informations $J(X),J(Y)$:
\begin{equation}\label{eq:classical-stam}
  \frac{1}{J(X+Y)}\;\ge\;\frac{1}{J(X)}+\frac{1}{J(Y)}.
\end{equation}
This is a cornerstone of information theory, closely related
to the entropy power inequality~\cite{blachman} and the
Cram\'er--Rao bound (see~\cite{dembo} for a survey).

Marcus, Spielman, and Srivastava~\cite{MSS2,MSS1}
introduced the \emph{finite free additive convolution}
$\fp$ on monic real-rooted polynomials of degree~$n$,
a finite-dimensional analogue of free additive
convolution in the sense
of~\cite{voiculescu}.
A natural question is whether the Stam
inequality~\eqref{eq:classical-stam} has a polynomial
analogue.
Define the \emph{finite free Fisher information} of
$r\in\PnR$ with simple roots
$\lambda_1<\cdots<\lambda_n$ by
$\Phi_n(r):=\sum_{i=1}^n
(\sum_{j\ne i}(\lambda_i-\lambda_j)^{-1})^{2}$.
The \emph{finite free Stam inequality}, conjectured by
D.~Shlyakhtenko and proved by Garza~Vargas, Srivastava,
and Stier~\cite{GVS}, states that for all $p,q\in\PnR$:
\begin{equation}\label{eq:stam-poly}
  \frac{1}{\Phi_n(p\fp q)}\;\ge\;
  \frac{1}{\Phi_n(p)}+\frac{1}{\Phi_n(q)}.
\end{equation}
Preserving real-rootedness under $\fp$ is guaranteed
by~\cite{MSS2}; see~\cite{bb} for the connection to
linear operators preserving stability.

\noindent\textbf{Contributions.}
\begin{enumerate}[nosep,label=(\roman*)]
\item Full proof of the Stam inequality for all $n\ge 2$
  (Section~\ref{sec:full-proof}), following~\cite{GVS}.
\item Structural identities
  (Section~\ref{sec:identities}):
  $\Phi_n=2\mathcal{R}=\tr(L)$, score identities,
  variance additivity, Bezoutian and Laplacian formulations.
\item The Hermite heat equation
  and root ODE (Section~\ref{sec:prelim});
  independent proofs for $n\le 3$ and
  Gaussian-input Stam at all~$n$
  (Section~\ref{sec:proved}).
\end{enumerate}

\smallskip
\noindent\textbf{Roadmap.}
Section~\ref{sec:prelim} establishes notation
and the key foundational results:
the heat equation for the Hermite flow
(Theorem~\ref{thm:hermite-heat}) and the
root ODE (Corollary~\ref{cor:root-ode}).
Section~\ref{sec:identities} develops
the Fisher--repulsion, Laplacian, and Bezoutian
identities.
Section~\ref{sec:proved} gives independent proofs
for small cases ($n\le 3$, Gaussian input).
Section~\ref{sec:full-proof} presents the full
proof for all~$n$.
Section~\ref{sec:discussion} concludes.


% ═══════════════════════════════════════════════════════════
\section{Preliminaries}\label{sec:prelim}
% ═══════════════════════════════════════════════════════════

\begin{definition}[MSS convolution~\cite{MSS2}]
\label{def:mss}
For $p(x)=\sum_{k=0}^n a_k x^{n-k}$ and $q(x)=\sum_{k=0}^n b_k x^{n-k}$ with $a_0=b_0=1$,
the \emph{finite free additive convolution} $r=p\fp q$ is defined by $r(x)=\sum_{k=0}^n c_k x^{n-k}$ with
\[
  c_k=\sum_{i+j=k}\frac{(n-i)!\,(n-j)!}{n!\,(n-k)!}\,a_i b_j.
\]
By~\cite{MSS2}, $\fp$ preserves $\PnR$.
\end{definition}

\begin{definition}[$K$-transform and log-cumulants]
\label{def:K}
Define
\[
  \kappa_k(r):=\frac{(n-k)!\,c_k(r)}{n!},
  \qquad
  K_r(z):=\sum_{k=0}^{n}\kappa_k(r)\,z^k.
\]
Then
\[
  K_{p\fp q}(z)=K_p(z)\cdot K_q(z)\pmod{z^{n+1}}.
\]
The \emph{log-cumulants} $\ell_k(r):=[z^k]\log K_r(z)$
satisfy $\ell_1=\kappa_1$ and
$\ell_k=\kappa_k-\tfrac{1}{k}\sum_{j=1}^{k-1}j\,\ell_j\kappa_{k-j}$
for $k\ge 2$.
They are \textbf{additive}:
$\ell_k(p\fp q)=\ell_k(p)+\ell_k(q)$ for all~$k$.
\end{definition}

\begin{definition}[Scores and Fisher information]
\label{def:scores}
For $r\in\PnR$ with simple roots
$\lambda_1<\cdots<\lambda_n$, let
$V_i(r):=\sum_{j\ne i}(\lambda_i-\lambda_j)^{-1}$
(the \emph{score vector} $V=(V_1,\ldots,V_n)$),
$\Phi_n(r):=\sum_i V_i^2$ (the \emph{Fisher information}),
$\mathcal{R}(r):=\sum_{i<j}(\lambda_i-\lambda_j)^{-2}$
(the \emph{repulsion energy}),
and $\mathcal{S}(r):=\sum_{i<j}(V_i-V_j)^2/(\lambda_i-\lambda_j)^2$
(the \emph{score-gradient energy}).
If $r$ has a repeated root, set $\Phi_n(r)=\infty$.
\end{definition}

\begin{definition}[Graph Laplacian]\label{def:laplacian}
The graph Laplacian of~$r$ is $L\in\R^{n\times n}$ with
$L_{ij}=-(\lambda_i-\lambda_j)^{-2}$ for $i\ne j$ and
$L_{ii}=\sum_{k\ne i}(\lambda_i-\lambda_k)^{-2}$.
We have $L\mathbf{1}=0$, $L\succeq 0$, $\rank L=n-1$.
Equivalently, $L=-\tfrac{1}{2}\mathrm{Hess}_\lambda(\log\disc(r))$.
\end{definition}

\begin{definition}[Variance and Gaussian polynomials]
\label{def:var}
For $r\in\PnR$: $\mu(r):=n^{-1}\sum_i\lambda_i$,
$\sigma^2(r):=n^{-1}\sum_i(\lambda_i-\mu)^2$.
Both are additive under~$\fp$
(Lemma~\ref{lem:var-add}).
The \emph{additive variance parameter}
$u:=\sigma^2/(2(n-1))$ satisfies
$u(p\fp q)=u(p)+u(q)$.
The \emph{finite Gaussian} $g_t\in\PnR$ has
$\sigma^2(g_t)=t$ and $\ell_k(g_t)=0$ for $k\ge 3$.
The Hermite semigroup satisfies $g_s\fp g_t=g_{s+t}$.
\end{definition}

\begin{definition}[Normalised cumulant ratios]
\label{def:tau}
For centred $r\in\PnR$ with $u:=-\ell_2(r)>0$, define
$\tau_k(r):=\ell_k(r)/u(r)^{k/2}$ for $k\ge 3$.
\end{definition}

\begin{lemma}[Normalisation identities]\label{lem:normalisation}
For centred $r\in\PnR$ (i.e., $\mu(r)=0$), the parameters
$\kappa_2$, $\ell_2$, $u$, and $\sigma^2$ are related by:
\begin{equation}\label{eq:norm-chain}
  \ell_2\;=\;\kappa_2\;=\;\frac{(n-2)!\,a_2}{n!}
  \;=\;\frac{a_2}{n(n-1)},
  \qquad
  u\;:=\;-\ell_2\;>\;0,
  \qquad
  \sigma^2\;=\;2(n-1)\,u.
\end{equation}
Here $a_2$ is the coefficient of $x^{n-2}$ in~$r$ (so $a_2<0$
for centred real-rooted $r$ with $n\ge 2$).
\end{lemma}

\begin{proof}
From Definition~\ref{def:K}: $\kappa_2=(n-2)!\,a_2/n!$.
The log-cumulant recurrence (Definition~\ref{def:K}) gives
$\ell_2=\kappa_2-\frac{1}{2}\kappa_1^2=\kappa_2$
when $r$ is centred ($\kappa_1=\ell_1=0$).
From the variance formula with $a_1=0$:
$\sigma^2=-2a_2/n=-2n(n-1)\ell_2/n=2(n-1)(-\ell_2)=2(n-1)u$.
All three parameters are additive under $\fp$
because $\ell_2$ is additive (Definition~\ref{def:K}).
\end{proof}


% ═══════════════════════════════════════════════════════════
\section{Structural identities}\label{sec:identities}
% ═══════════════════════════════════════════════════════════

We collect the main identities connecting
$\Phi_n$ to spectral quantities.  Throughout
this section, $r\in\PnR$ has simple roots.

\begin{theorem}[Fisher--repulsion identity]
\label{thm:fisher-rep}
$\Phi_n(r)=2\,\mathcal{R}(r)$.
\end{theorem}

\begin{proof}
Expand $\Phi_n=\sum_i V_i^2=\sum_i\sum_{j\ne i}\sum_{k\ne i}(\lambda_i-\lambda_j)^{-1}(\lambda_i-\lambda_k)^{-1}$.
The diagonal terms ($j=k$) sum to $2\sum_{i<j}(\lambda_i-\lambda_j)^{-2}=2\mathcal{R}$.
The cross-terms ($j\ne k$, both $\ne i$) group into triples $\{a,b,c\}$, each contributing
$(a-b)^{-1}(a-c)^{-1}+(b-a)^{-1}(b-c)^{-1}+(c-a)^{-1}(c-b)^{-1}=0$
by the partial-fraction identity.
\end{proof}

\begin{theorem}[Fisher--Laplacian identities]
\label{thm:fisher-lap}
\begin{enumerate}[label=\textup{(\alph*)},nosep]
  \item $\Phi_n=\tr(L)$.
  \item $V=L\lambda$ \textup{(Euler identity)}.
  \item $\lambda^T L\lambda=\binom{n}{2}$.
  \item $\Phi_n=\|L\lambda\|^2=\lambda^T L^2\lambda$.
\end{enumerate}
\end{theorem}

\begin{proof}
(a) $\tr(L)=\sum_i\sum_{k\ne i}(\lambda_i-\lambda_k)^{-2}=2\sum_{i<j}(\lambda_i-\lambda_j)^{-2}=2\mathcal{R}=\Phi_n$.

(b) $(L\lambda)_i=\sum_{j\ne i}(\lambda_i-\lambda_j)/(\lambda_i-\lambda_j)^2=\sum_{j\ne i}(\lambda_i-\lambda_j)^{-1}=V_i$.

(c) $\lambda^T L\lambda=V\cdot\lambda=\sum_i\lambda_i V_i=\binom{n}{2}$ by the Euler identity for $\disc$ (degree $n(n-1)$ homogeneous).

(d) Immediate from $V=L\lambda$.
\end{proof}

\begin{lemma}[Score identities]\label{lem:score-ids}
\begin{enumerate}[label=\textup{(\roman*)},nosep]
  \item $\sum_i V_i=0$.
  \item $\sum_i(\lambda_i-\mu)V_i=\binom{n}{2}$.
  \item $\Phi_n=\sum_{i<j}(V_i-V_j)/(\lambda_i-\lambda_j)$.
  \item $V_i=r''(\lambda_i)/(2r'(\lambda_i))$.
\end{enumerate}
\end{lemma}

\begin{proof}
(i) $\sum_i V_i=\sum_{i\ne j}(\lambda_i-\lambda_j)^{-1}=0$ (antisymmetric).

(ii) $\sum_i\lambda_i V_i=\sum_{i\ne j}\lambda_i/(\lambda_i-\lambda_j)
=\sum_{i\ne j}\bigl[1+\lambda_j/(\lambda_i-\lambda_j)\bigr]
=n(n-1)+\sum_{i\ne j}\lambda_j/(\lambda_i-\lambda_j)$.
Using $\sum_{i\ne j}\lambda_j/(\lambda_i-\lambda_j)
=-\sum_{i\ne j}\lambda_i/(\lambda_j-\lambda_i)
=-\sum_i\lambda_i V_i$,
we get $2\sum_i\lambda_i V_i=n(n-1)$,
so $\sum_i\lambda_i V_i=\binom{n}{2}$.
By (i), subtracting $\mu\sum V_i=0$ gives (ii).

(iii) Expand the right-hand side:
\begin{align*}
  \sum_{i<j}\frac{V_i-V_j}{\lambda_i-\lambda_j}
  &=\sum_{i<j}\frac{1}{\lambda_i-\lambda_j}
    \Bigl(\sum_{k\ne i}\frac{1}{\lambda_i-\lambda_k}
    -\sum_{k\ne j}\frac{1}{\lambda_j-\lambda_k}\Bigr)\\
  &=\sum_{i<j}\sum_{k\ne i}\frac{1}{(\lambda_i-\lambda_j)(\lambda_i-\lambda_k)}
    -\sum_{i<j}\sum_{k\ne j}\frac{1}{(\lambda_i-\lambda_j)(\lambda_j-\lambda_k)}.
\end{align*}
Relabelling $i\leftrightarrow j$ in the second sum and combining yields
$2\sum_{i<j}\sum_{k\ne i}1/((\lambda_i-\lambda_j)(\lambda_i-\lambda_k))$.
Separating diagonal ($k=j$) from cross ($k\ne i,j$) terms:
the diagonal gives $\sum_i\sum_{j\ne i}(\lambda_i-\lambda_j)^{-2}=\Phi_n$;
the cross-terms group into triples $\{i,j,k\}$, each contributing
$\sum_{\mathrm{cyc}}1/((a-b)(a-c))=0$
by the same partial-fraction identity as in Theorem~\ref{thm:fisher-rep}.

(iv) Since $r'(\lambda_i)=\prod_{j\ne i}(\lambda_i-\lambda_j)$, we have $V_i=r''(\lambda_i)/(2r'(\lambda_i))$.
\end{proof}

\begin{theorem}[Fisher--variance inequality]
\label{thm:fisher-var}
$\Phi_n(r)\cdot\sigma^2(r)\ge n(n-1)^2/4$.
\end{theorem}

\begin{proof}
Cauchy--Schwarz on
$\sum_i(\lambda_i-\mu)V_i=\binom{n}{2}$
with $\sum V_i=0$:
$\bigl|\sum(\lambda_i-\mu)V_i\bigr|^2
\le\bigl(\sum(\lambda_i-\mu)^2\bigr)\bigl(\sum V_i^2\bigr)
=n\sigma^2\cdot\Phi_n$.
Hence $n\sigma^2\cdot\Phi_n\ge\binom{n}{2}^2=n^2(n-1)^2/4$.
\end{proof}

\begin{theorem}[Score-gradient inequality]
\label{thm:sgi}
$\mathcal{S}(r)\cdot\sigma^2(r)\ge(n-1)\Phi_n(r)/2$.
\end{theorem}

\begin{proof}
Write $\lambda_c:=\lambda-\mu\mathbf{1}$ for the centred root vector.
Since $L\mathbf{1}=0$, $V=L\lambda=L\lambda_c$.
The Cauchy--Schwarz inequality for the positive semi-definite form
$\langle u,v\rangle_L:=u^TLv$ gives
$(\lambda_c^T L^2\lambda_c)^2
\le(\lambda_c^T L\lambda_c)(\lambda_c^T L^3\lambda_c)$,
i.e., $\Phi_n^2\le\binom{n}{2}\cdot\mathcal{S}$.
Combining with the Fisher--variance inequality
(Theorem~\ref{thm:fisher-var}):
\[
  \mathcal{S}\,\sigma^2
  \;\ge\;\frac{\Phi_n^2}{\binom{n}{2}}\,\sigma^2
  \;=\;\frac{\Phi_n\,\sigma^2}{\binom{n}{2}}\cdot\Phi_n
  \;\ge\;\frac{n(n-1)^2/4}{n(n-1)/2}\cdot\Phi_n
  \;=\;\frac{(n-1)\,\Phi_n}{2}.\qedhere
\]
\end{proof}

\begin{theorem}[Bezoutian representation]
\label{thm:bez}
$\Phi_n(r)=\sum_{i=1}^n r''(\lambda_i)^2
/(4\,r'(\lambda_i)^2)
=\|r''/2\|^2_{\Bez(r,r')}$.
\end{theorem}

\begin{proof}
The Bezoutian matrix $\Bez(r,r')$ is the unique symmetric
$B\in\R^{n\times n}$ satisfying
$\sum_{i,j}B_{ij}x^{n-1-i}y^{n-1-j}
=\bigl(r(x)r'(y)-r'(x)r(y)\bigr)/(x-y)$.
The associated inner product is diagonal in the Lagrange basis
$\{L_i(x)=\prod_{j\ne i}(x-\lambda_j)/\prod_{j\ne i}(\lambda_i-\lambda_j)\}$:
$\langle f,g\rangle_{\Bez}=\sum_i f(\lambda_i)g(\lambda_i)/r'(\lambda_i)^2$
(see~\cite{fiedler} for the diagonalisation).
Since $V_i=r''(\lambda_i)/(2r'(\lambda_i))$
(Lemma~\ref{lem:score-ids}(iv)), we get
$\Phi_n=\sum V_i^2=\|r''/2\|^2_{\Bez}$.
\end{proof}

\begin{lemma}[Variance additivity]\label{lem:var-add}
$\sigma^2(p\fp q)=\sigma^2(p)+\sigma^2(q)$.
\end{lemma}

\begin{proof}
From the MSS coefficient formula (Definition~\ref{def:mss}):
$c_1=a_1+b_1$, $c_2=a_2+b_2+\frac{n-1}{n}a_1 b_1$.
Using $\sigma^2=\frac{(n-1)a_1^2-2na_2}{n^2}$:
\begin{align*}
  \sigma^2(p\fp q)
  &=\frac{(n\!-\!1)(a_1\!+\!b_1)^2-2n(a_2\!+\!b_2\!+\!\tfrac{n-1}{n}a_1 b_1)}{n^2}\\
  &=\frac{(n\!-\!1)a_1^2\!-\!2na_2}{n^2}
   +\frac{(n\!-\!1)b_1^2\!-\!2nb_2}{n^2}
   +\frac{2(n\!-\!1)a_1b_1\!-\!2(n\!-\!1)a_1b_1}{n^2}\\
  &=\sigma^2(p)+\sigma^2(q).\qedhere
\end{align*}
\end{proof}

\begin{lemma}[Derivative compatibility]\label{lem:deriv}
$(p\fp q)'/n=(p'/n)\boxplus_{n-1}(q'/n)$.
\end{lemma}

\begin{proof}
The monic degree-$(n-1)$ polynomial $p'/n$ has
coefficients $\tilde a_k=(n-k)a_k/n$.
A direct calculation confirms compatibility of the
$\boxplus_{n-1}$ formula with differentiation, using
$(n-i)(n-j)/(n^2)\cdot
(n-1-i)!(n-1-j)!/((n-1)!(n-1-k)!)
=(n-k)/n\cdot(n-i)!(n-j)!/(n!(n-k)!)$ for $i+j=k$.
\end{proof}

\begin{theorem}[Hermite heat equation]\label{thm:hermite-heat}
The Hermite flow $r_t:=r\fp g_t$ satisfies the
scaled backwards heat equation
\begin{equation}\label{eq:heat}
  \partial_t r_t(x)
  \;=\;-\frac{1}{2(n-1)}\,r_t''(x),
\end{equation}
with initial condition $r_0=r$.
\end{theorem}

\begin{proof}
The $K$-transform of the finite Gaussian is
$K_{g_t}(z)=\exp\bigl(-\tfrac{t}{2(n-1)}z^2\bigr)
\bmod z^{n+1}$,
since $\ell_2(g_t)=-t/(2(n-1))$ and
$\ell_k(g_t)=0$ for $k\ge 3$
(Definition~\ref{def:var}).
Hence
$K_{r_t}(z)=K_r(z)\,K_{g_t}(z)$
gives
$\partial_t K_{r_t}(z)
=-\tfrac{z^2}{2(n-1)}\,K_{r_t}(z)
\bmod z^{n+1}$.
Expanding $K_{r_t}(z)=\sum_{k=0}^n\kappa_k(t)\,z^k$
yields the coefficient ODE
\[
  \kappa_k'(t)
  =-\frac{1}{2(n-1)}\,\kappa_{k-2}(t)
  \quad(k\ge 2),
  \qquad
  \kappa_0'=\kappa_1'=0.
\]
Translating to the standard coefficients
$c_k(t)=n!\,\kappa_k(t)/(n-k)!$:
\[
  c_k'(t)
  =-\frac{(n-k+2)(n-k+1)}{2(n-1)}\,c_{k-2}(t)
  \qquad(k\ge 2).
\]
Now compute
$\partial_t r_t(x)
=\sum_{k=0}^n c_k'(t)\,x^{n-k}$
and observe that
$r_t''(x)
=\sum_{j=0}^{n-2}(n{-}j)(n{-}j{-}1)\,c_j(t)\,x^{n-j-2}$;
substituting $k=j+2$ identifies the two sums,
giving~\eqref{eq:heat}.
\end{proof}

\begin{remark}[Context]\label{rem:heat-context}
The heat equation~\eqref{eq:heat} is exact
for polynomials of degree~$n$
(no truncation is involved, since all
coefficients lie in degrees $\le n$).
The root ODE~\eqref{eq:root-ode} is a deterministic
analogue of the Dyson Brownian motion
for eigenvalues of random Hermitian
matrices~\cite{dyson};
here it arises purely algebraically
from the $K$-transform.
\end{remark}

\begin{corollary}[Root ODE]\label{cor:root-ode}
Along the Hermite flow $r_t=r\fp g_t$, each simple
root $\lambda_i(t)$ of $r_t$ satisfies
\begin{equation}\label{eq:root-ode}
  \dot\lambda_i(t)
  \;=\;\frac{V_i(t)}{n-1}
  \;=\;\frac{1}{n-1}\sum_{j\ne i}
  \frac{1}{\lambda_i(t)-\lambda_j(t)}.
\end{equation}
\end{corollary}

\begin{proof}
Differentiating $r_t(\lambda_i(t))=0$ in~$t$ and
applying Theorem~\ref{thm:hermite-heat}:
\[
  0=\partial_t r_t(\lambda_i)
  +r_t'(\lambda_i)\,\dot\lambda_i
  =-\frac{r_t''(\lambda_i)}{2(n-1)}
  +r_t'(\lambda_i)\,\dot\lambda_i,
\]
so $\dot\lambda_i
=r_t''(\lambda_i)/(2(n-1)\,r_t'(\lambda_i))
=V_i/(n-1)$
by Lemma~\ref{lem:score-ids}\textup{(iv)}.
\end{proof}

\begin{theorem}[De Bruijn identity]\label{thm:debruijn}
Along the Hermite flow $r_t=r\fp g_t$:
$\frac{d}{dt}\log|\disc(r_t)|
=\frac{2}{n-1}\Phi_n(r_t)$.
\end{theorem}

\begin{proof}
By Corollary~\ref{cor:root-ode},
$\dot\lambda_i=V_i/(n-1)$.
Since $\disc(r)=\prod_{i<j}(\lambda_i-\lambda_j)^2$,
we have $\partial_{\lambda_i}\log\disc=2V_i$.
Therefore
$\frac{d}{dt}\log\disc
=\sum_i 2V_i\cdot V_i/(n-1)
=2\Phi_n/(n-1)$.
\end{proof}


% ═══════════════════════════════════════════════════════════
\section{Proved cases of the Stam inequality}
\label{sec:proved}
% ═══════════════════════════════════════════════════════════

\subsection{The case \texorpdfstring{$n=2$}{n=2}: equality}
\label{ssec:n2}

For $n=2$: $\Phi_2(r)=2/(\lambda_1-\lambda_2)^2
=1/(2\sigma^2)$, so $1/\Phi_2=2\sigma^2$.
The Stam inequality reduces to variance additivity
(Lemma~\ref{lem:var-add}), with equality.

\subsection{The case \texorpdfstring{$n=3$}{n=3}: SOS proof}
\label{ssec:n3}

\begin{theorem}[Stam for \texorpdfstring{$n=3$}{n=3}]
\label{thm:n3-stam}
For centred $p,q\in\PnR[3]$ with $u_p,u_q>0$,
let $r=p\fp q$, $w=u_p/(u_p+u_q)$,
$\alpha:=\ell_3(p)/u_p$, $\beta:=\ell_3(q)/u_q$.
Then
\begin{equation}\label{eq:D3}
  D_3:=\frac{1}{\Phi_3(r)}-\frac{1}{\Phi_3(p)}
  -\frac{1}{\Phi_3(q)}
  =\frac{3}{2}\bigl[(1-w)\alpha^2
  +w(1-w)(\alpha-\beta)^2+w\beta^2\bigr]\ge 0.
\end{equation}
Equality holds \textup{(}for $w\in(0,1)$\textup{)}
iff $\ell_3(p)=\ell_3(q)=0$.
\end{theorem}

\begin{proof}
\emph{Step~1: Log-cumulants for the depressed cubic.}
Let $r(x)=x^3+e_2 x+e_3$ be a centred monic cubic.
The $K$-transform coefficients
(Definition~\ref{def:K}) are
$\kappa_0=1$, $\kappa_1=0$,
$\kappa_2=\frac{1!\cdot e_2}{3!}=\frac{e_2}{6}$,
$\kappa_3=\frac{0!\cdot e_3}{3!}=\frac{e_3}{6}$,
so $K_r(z)=1+\frac{e_2}{6}z^2+\frac{e_3}{6}z^3$.
Since $\log(1+x)=x-\tfrac{x^2}{2}+\cdots$
and $K_r$ has no $z^1$ term:
$\ell_2=[z^2]\log K_r=\kappa_2=e_2/6$,
$\ell_3=[z^3]\log K_r=\kappa_3=e_3/6$.
(The cross-term $\kappa_2^2 z^4/2$
contributes to $[z^4]\log K_r$,
not to~$[z^3]$.)
Hence $u:=-\ell_2=-e_2/6>0$ and
$e_2=-6u$, $e_3=6\ell_3$.

\emph{Step~2: $\Phi_3$ via
the repulsion energy.}
Denote the roots $\lambda_1<\lambda_2<\lambda_3$
and the three squared gaps
$D_{ij}:=(\lambda_i-\lambda_j)^2$.
Since $r'(\lambda_i)
=\prod_{j\ne i}(\lambda_i-\lambda_j)$,
we have
$r'(\lambda_i)^2
= \prod_{j\ne i}D_{ij}$.
Hence the sum of products of two
squared gaps equals
$\sum_i r'(\lambda_i)^2$.
For $r'(x)=3x^2+e_2$ and
the Newton power sums
$p_2=-2e_2$, $p_4=2e_2^2$:
\[
  \sum_{i=1}^3 r'(\lambda_i)^2
  = \sum_i (3\lambda_i^2+e_2)^2
  = 9\,p_4+6\,e_2\,p_2+3\,e_2^2
  = 18e_2^2-12e_2^2+3e_2^2
  = 9e_2^2.
\]
The discriminant is
$\Delta_3=\prod_{i<j}D_{ij}
  =-4e_2^3-27e_3^2
  =864\,u^3-972\,\ell_3^2$,
and the repulsion energy decomposes as
\[
  \mathcal{R}
  =\sum_{i<j}\frac{1}{D_{ij}}
  =\frac{\sum_i r'(\lambda_i)^2}{\Delta_3}
  =\frac{9e_2^2}{\Delta_3}.
\]
By Theorem~\ref{thm:fisher-rep},
$\Phi_3=2\mathcal{R}
  =18e_2^2/\Delta_3
  =648\,u^2/(864\,u^3-972\,\ell_3^2)$.

\emph{Step~3: Closed-form reciprocal.}
\begin{equation}\label{eq:inv-phi3}
  \frac{1}{\Phi_3(r)}
  =\frac{864\,u^3-972\,\ell_3^2}{648\,u^2}
  =\frac{4u}{3}-\frac{3\ell_3^2}{2u^2},
\end{equation}
where the last equality uses
$864/648=4/3$ and $972/648=3/2$.

\emph{Step~4: Defect computation.}
Since $u$ and $\ell_3$ are additive
under~$\fp$, set $u_r=u_p+u_q$ and
$\ell_{3,r}=\ell_{3,p}+\ell_{3,q}$.
With $\alpha=\ell_{3,p}/u_p$ and
$\beta=\ell_{3,q}/u_q$:
\begin{align*}
  D_3 &= \frac{4u_r}{3}
         -\frac{3\ell_{3,r}^2}{2u_r^2}
         -\frac{4u_p}{3}
         +\frac{3\ell_{3,p}^2}{2u_p^2}
         -\frac{4u_q}{3}
         +\frac{3\ell_{3,q}^2}{2u_q^2}\\[4pt]
      &= \frac{3}{2}\biggl[
           \alpha^2+\beta^2
           -\frac{(u_p\alpha+u_q\beta)^2}
                 {(u_p+u_q)^2}
         \biggr],
\end{align*}
since the linear terms
$\tfrac{4}{3}(u_r{-}u_p{-}u_q)=0$
cancel by additivity.
Writing $w=u_p/(u_p+u_q)$:
\begin{align*}
  D_3 &= \frac{3}{2}\bigl[
    \alpha^2+\beta^2
    -w^2\alpha^2
    -2w(1{-}w)\alpha\beta
    -(1{-}w)^2\beta^2
  \bigr]\\[4pt]
  &= \frac{3}{2}\bigl[
    (1{-}w^2)\alpha^2
    -2w(1{-}w)\alpha\beta
    +w(2{-}w)\beta^2
  \bigr].
\end{align*}
Since $(1{-}w)(1{+}w)=1{-}w^2$
and $w(2{-}w)=w(1{-}w)+w$:
\[
  D_3 = \frac{3}{2}\bigl[
    (1{-}w)\alpha^2
    +w(1{-}w)(\alpha{-}\beta)^2
    +w\beta^2
  \bigr]\ge 0.
\]
Each of the three summands is
manifestly non-negative;
for $w\in(0,1)$, all vanish iff
$\alpha=\beta=0$,
i.e.\ $\ell_3(p)=\ell_3(q)=0$.
\end{proof}

\begin{remark}\label{rem:n3-mech}
The structure of $1/\Phi_3$ is
$1/\Phi_3=A(u)+Q(\ell_3/u)$ where
$A(u)=4u/3$ is additive under~$\fp$ and
$Q(\cdot)=-\tfrac{3}{2}(\cdot)^2$ is concave.
The Stam defect therefore reduces to
the concavity defect of a quadratic
composed with a weighted-linear mixing law.
The Hessian of $1/\Phi_3$ in
$(u,\ell_3)$-coordinates is \textbf{not}
negative semi-definite, so the result does
not follow from global concavity of
$1/\Phi_3$ as a function of both variables.
\end{remark}

\subsection{Gaussian-input Stam for all \texorpdfstring{$n$}{n}}

\begin{theorem}[Gaussian-input Stam inequality]
\label{thm:stam-gauss}
For all $r\in\PnR$ and $t>0$:
$1/\Phi_n(r\fp g_t)\ge 1/\Phi_n(r)+1/\Phi_n(g_t)$.
Equality holds on the Hermite manifold.
\end{theorem}

\begin{proof}
\emph{Step~1: Gaussian Fisher information.}
The finite Gaussian $g_t$ has roots at the
$n$~zeros of the probabilist Hermite polynomial $\mathrm{He}_n$
scaled so that $\sigma^2(g_t)=t$ and
$\ell_k(g_t)=0$ for $k\ge 3$.
The Hermite differential equation
$\mathrm{He}_n''(x)-x\,\mathrm{He}_n'(x)+n\,\mathrm{He}_n(x)=0$
at a root $\lambda_i$ (where $\mathrm{He}_n(\lambda_i)=0$) gives
$\mathrm{He}_n''(\lambda_i)=\lambda_i\,\mathrm{He}_n'(\lambda_i)$,
hence
$V_i=\mathrm{He}_n''(\lambda_i)/(2\mathrm{He}_n'(\lambda_i))
=\lambda_i/2$.
After scaling, $V_i(g_t)=c\,(\lambda_i-\mu)$
for a constant $c$ depending only on $n$ and $t$,
so equality holds in the Fisher--variance inequality
(Theorem~\ref{thm:fisher-var}):
$\Phi_n(g_t)=n(n-1)^2/(4t)$
and $1/\Phi_n(g_t)=4t/(n(n-1)^2)$.

\emph{Step~2: Root ODE.}
Let $r_t:=r\fp g_t$ with roots
$\lambda_1(t)<\cdots<\lambda_n(t)$.
By Corollary~\ref{cor:root-ode},
each root satisfies
$\dot\lambda_i = V_i/(n-1)$,
where $V_i=V_i(r_t)$.

\emph{Step~3: Score ODE.}
Differentiating
$V_i=\sum_{j\ne i}(\lambda_i-\lambda_j)^{-1}$
with respect to~$t$:
\[
  \dot V_i
  = -\sum_{j\ne i}
    \frac{\dot\lambda_i-\dot\lambda_j}
         {(\lambda_i-\lambda_j)^2}
  = -\frac{1}{n-1}\sum_{j\ne i}
    \frac{V_i-V_j}{(\lambda_i-\lambda_j)^2}
  = -\frac{(LV)_i}{n-1},
\]
where $L$ is the graph Laplacian
(Definition~\ref{def:laplacian}).

\emph{Step~4: Fisher dissipation.}
Define the \emph{score-gradient energy}
$\mathcal{S}:=V^T LV
  =\sum_{i<j}(V_i-V_j)^2
   /(\lambda_i-\lambda_j)^2$.
Then
$\frac{d}{dt}\Phi_n(r_t)
  = 2\sum_i V_i\dot V_i
  = -\frac{2}{n-1}V^T L V
  = -\frac{2}{n-1}\mathcal{S}$.
Since $L\succeq 0$, we have
$\mathcal{S}\ge 0$, hence
$\Phi_n'(r_t)\le 0$.

\emph{Step~5: Lower bound on $\mathcal{S}/\Phi_n^2$.}
Since $L\mathbf{1}=0$, $V=L\lambda=L\lambda_c$
where $\lambda_c:=\lambda-\mu\mathbf{1}$.
Hence $\Phi_n=V^T V=\lambda_c^T L^2\lambda_c$
and $\mathcal{S}=V^T LV=\lambda_c^T L^3\lambda_c$.
The Cauchy--Schwarz inequality for the
semi-definite form $\langle u,v\rangle_L=u^T Lv$ gives
$(\lambda_c^T L^2\lambda_c)^2
\le(\lambda_c^T L\lambda_c)
   (\lambda_c^T L^3\lambda_c)
=\binom{n}{2}\,\mathcal{S}$,
hence $\Phi_n^2\le\binom{n}{2}\,\mathcal{S}$,
i.e.,
\[
  \frac{\mathcal{S}}{\Phi_n^2}
  \;\ge\;\frac{1}{\binom{n}{2}}
  \;=\;\frac{2}{n(n-1)}.
\]

\emph{Step~6: Integration.}
Since
$\bigl(1/\Phi_n\bigr)'
  = -\Phi_n'/(\Phi_n^2)
  = \frac{2}{n-1}\cdot
    \frac{\mathcal{S}}{\Phi_n^2}
  \ge \frac{2}{n-1}\cdot
      \frac{2}{n(n-1)}
  = \frac{4}{n(n-1)^2}$,
integrating from $0$ to $t$:
\[
  \frac{1}{\Phi_n(r_t)}
  -\frac{1}{\Phi_n(r)}
  \ge \frac{4t}{n(n-1)^2}
  = \frac{1}{\Phi_n(g_t)}.
\]

\emph{Equality saturation.}
When $r=g_s$ is itself a finite Gaussian,
$r_t=g_s\fp g_t=g_{s+t}$,
and the scores satisfy
$V_i(g_u)=c(\lambda_i-\mu)$ for every $u>0$
(Step~1).
Then $\mathcal{S}/\Phi_n^2=1/\binom{n}{2}$ exactly,
so $(1/\Phi_n)'=4/(n(n-1)^2)$ for all~$t>0$,
and the integrated bound is attained with equality.
\end{proof}


% ═══════════════════════════════════════════════════════════
\section{Proof of the Stam inequality for all \texorpdfstring{$n$}{n}}
\label{sec:full-proof}
% ═══════════════════════════════════════════════════════════

We now present the full proof of the finite free Stam
inequality, following Garza~Vargas, Srivastava, and
Stier~\cite{GVS}.
The argument uses three ingredients:
(1)~a relation between score vectors mediated
by the Jacobian of the root map;
(2)~a contraction inequality for this Jacobian,
derived from variance additivity and a convexity
result on hyperbolic polynomials;
(3)~Blachman's classical argument~\cite{blachman}.

\subsection{The root map and its Jacobian}

\begin{definition}[Root map]\label{def:root-map}
For $p,q\in\PnR$ with root vectors
$\alpha=(\alpha_1\ge\cdots\ge\alpha_n)$ and
$\beta=(\beta_1\ge\cdots\ge\beta_n)$,
define $\Omega_{\fp}(\alpha,\beta)\in\R^n$ to be
the ordered root vector of $p\fp q$.
By~\cite{MSS2}, $\Omega_{\fp}$ is well-defined and
takes values in \emph{ordered} real vectors.
Write $\gamma=\Omega_{\fp}(\alpha,\beta)$.
Let $J_{\fp}$ denote the Jacobian
(the $n\times 2n$ matrix of partial derivatives)
of $\Omega_{\fp}$ at $(\alpha,\beta)$,
defined whenever the roots are simple.
\end{definition}

\begin{lemma}[Approximation by simple roots]
\label{lem:approx}
Let $T_\varepsilon:=(1-\varepsilon\,d/dx)^n$.
For any monic real-rooted $p$ of degree~$n$:
\begin{enumerate}[label=\textup{(\alph*)},nosep]
  \item $T_\varepsilon p$ is monic, real-rooted of
    degree~$n$, with simple roots
    \textup{(}Nuij~\cite{nuij}\textup{)}.
  \item $\Phi_n(T_\varepsilon p)\to\Phi_n(p)$
    as $\varepsilon\to 0$.
  \item $(T_\varepsilon p)\fp(T_\varepsilon q)
    =T_\varepsilon^2(p\fp q)$,
    since $\fp$ commutes with differential
    operators~\cite{MSS2}.
\end{enumerate}
Hence it suffices to prove
Theorem~\textup{\ref{thm:stam-full}}
when $p$, $q$, and $p\fp q$ all have
simple roots.
\end{lemma}

\subsection{Score vectors and the Jacobian}

\begin{proposition}[Score--Jacobian relation]
\label{prop:score-jac}
Let $p,q\in\PnR$ have simple roots with
root vectors $\alpha,\beta$ and
$\gamma=\Omega_{\fp}(\alpha,\beta)$.
For any $a,b\ge 0$:
\begin{equation}\label{eq:score-jac}
  J_{\fp}\begin{pmatrix}a\,V(\alpha)\\b\,V(\beta)\end{pmatrix}
  =(a+b)\,V(\gamma),
\end{equation}
where $V(\alpha)=\bigl(\sum_{j\ne i}(\alpha_i-\alpha_j)^{-1}\bigr)_{i=1}^n$
is the score vector.
\end{proposition}

\begin{proof}
For $t\ge 0$, define
$p_t:=\exp(-\tfrac{t}{2}\partial_x^2)\,p$
with ordered root vector $\alpha(t)$, and
similarly $q_t,r_t$ with $\beta(t),\gamma(t)$.
Since $\fp$ commutes with any differential
operator, $r_{(a+b)t}=p_{at}\fp q_{bt}$,
whence $\gamma((a+b)t)
=\Omega_{\fp}(\alpha(at),\beta(bt))$.
Differentiating at $t=0$ and using the chain rule:
\[
  (a+b)\,\gamma'(0)
  =J_{\fp}\begin{pmatrix}
    a\,\alpha'(0)\\b\,\beta'(0)
  \end{pmatrix}.
\]
By Lemma~\ref{lem:score-ids}\textup{(iv)},
$\alpha'_i(0)
=p''(\alpha_i)/(2p'(\alpha_i))
=V_i(\alpha)$
(the same computation as in
Corollary~\ref{cor:root-ode}).
\end{proof}

\subsection{Jacobian contraction}

\begin{proposition}[Jacobian contraction on $V$]
\label{prop:jac-contract}
Let $V:=\{(u,v)\in\R^n\times\R^n:
u^T\mathbf{1}=v^T\mathbf{1}=0\}$.
For all $(u,v)\in V$:
\begin{equation}\label{eq:jac-contract}
  \bigl\|J_{\fp}(u,v)\bigr\|^2
  \;\le\;\|u\|^2+\|v\|^2.
\end{equation}
\end{proposition}

The proof of this proposition requires two
intermediate results.

\begin{lemma}[Hessian identity]
\label{lem:hessian-id}
Write $H^{(i)}_{\fp}:=\mathrm{Hess}\,\Omega_{\fp,i}$
for the Hessian of the $i$-th coordinate function
of $\Omega_{\fp}$ at $(\alpha,\beta)$.
Then for all $w=(u,v)\in V$:
\begin{equation}\label{eq:hessian-id}
  w^T J_{\fp}J_{\fp}^T w
  \;=\;\|w\|^2
  -\sum_{i=1}^n\gamma_i\,w^T H^{(i)}_{\fp}\,w.
\end{equation}
\end{lemma}

\begin{proof}
Fix $w=(u,v)\in V$ and consider the paths
$\alpha(t):=\alpha+tu$, $\beta(t):=\beta+tv$,
$\gamma(t):=\Omega_{\fp}(\alpha(t),\beta(t))$.
By variance additivity (Lemma~\ref{lem:var-add})
applied to root vectors:
\[
  \frac{1}{n}\sum_i\gamma_i(t)^2
  -\Bigl(\frac{1}{n}\sum_i\gamma_i(t)\Bigr)^{\!2}
  =\mathrm{Var}(\alpha(t))+\mathrm{Var}(\beta(t)).
\]
Since $(u,v)\in V$, the means
$\frac{1}{n}\sum\alpha_i(t)$
and $\frac{1}{n}\sum\beta_i(t)$
(and hence $\frac{1}{n}\sum\gamma_i(t)$) are
constant in~$t$.
Differentiating the equation
$\frac{1}{n}\sum_i\gamma_i(t)^2
=\mathrm{Var}(\alpha(t))+\mathrm{Var}(\beta(t))
+\mathrm{const}$
twice at $t=0$:
\[
  \frac{1}{n}\sum_i\bigl[
    (D_w\gamma_i)^2+\gamma_i\,D_w^2\gamma_i
  \bigr]
  =\frac{1}{n}(\|u\|^2+\|v\|^2),
\]
since $\partial_t^2[\|\alpha+tu\|^2]=2\|u\|^2$
and similarly for~$v$.
Now $\sum_i(D_w\gamma_i)^2=w^T J_{\fp}J_{\fp}^T w$
and $D_w^2\gamma_i=w^T H^{(i)}_{\fp}w$,
giving~\eqref{eq:hessian-id}.
\end{proof}

\begin{lemma}[PSD Hessian sum]
\label{lem:psd-hess}
For any $c_1\ge\cdots\ge c_n$,
the matrix $\sum_{i=1}^n c_i\,H^{(i)}_{\fp}$
is positive semi-definite.
\end{lemma}

\begin{proof}
Consider the multivariate polynomial
\[
  f(x,a_1,\ldots,a_n,b_1,\ldots,b_n)
  :=\mathop{\mathbb{E}}_{\pi\in S_n}
  \prod_{i=1}^n(x-a_i-b_{\pi(i)}),
\]
which is homogeneous and, since $\fp$ preserves
real-rootedness, hyperbolic in the direction
$e_1=(1,0,\ldots,0)$.
Its roots in~$x$ (at fixed $a=\alpha$, $b=\beta$)
are the entries $\gamma_i$ of
$\Omega_{\fp}(\alpha,\beta)$.
By a theorem of Bauschke, G\"uler, Lewis,
and Sendov~\cite{BGLS},
the partial sums
$\sigma_k(a,b):=\sum_{i=1}^k\gamma_i(a,b)$
are convex functions; in particular,
$\mathrm{Hess}\,\sigma_k\succeq 0$.
Since the $c_i$ are decreasing,
$\sum c_i\gamma_i$ is a positive linear
combination of the~$\sigma_k$,
hence $\mathrm{Hess}(\sum c_i\gamma_i)
=\sum c_i H^{(i)}_{\fp}\succeq 0$.
\end{proof}

\begin{proof}[Proof of
Proposition~\textup{\ref{prop:jac-contract}}]
For $(u,v)\in V$, by
Lemma~\ref{lem:hessian-id}:
\[
  \|J_{\fp}(u,v)\|^2
  =\|u\|^2+\|v\|^2
  -\sum_{i=1}^n\gamma_i\,w^TH^{(i)}_{\fp}\,w.
\]
The roots $\gamma_1\ge\cdots\ge\gamma_n$ are ordered,
so Lemma~\ref{lem:psd-hess} (with $c_i=\gamma_i$)
gives $\sum\gamma_i\,w^TH^{(i)}_{\fp}\,w\ge 0$.
\end{proof}

\subsection{Main theorem}

\begin{theorem}[Finite free Stam inequality]
\label{thm:stam-full}
For all $n\ge 2$ and $p,q\in\PnR$:
\begin{equation}\label{eq:stam-full}
  \frac{1}{\Phi_n(p\fp q)}
  \;\ge\;\frac{1}{\Phi_n(p)}+\frac{1}{\Phi_n(q)}.
\end{equation}
\end{theorem}

\begin{proof}
By Lemma~\ref{lem:approx}, we may assume
$p$, $q$, and $r:=p\fp q$ all have simple roots.
Write $\alpha,\beta,\gamma$ for their ordered
root vectors.

\emph{Step~1: Score vectors are in~$V$.}
By Lemma~\ref{lem:score-ids}\textup{(i)},
$\sum_i V_i(\alpha)=0$ and
$\sum_i V_i(\beta)=0$,
so $(V(\alpha),V(\beta))\in V$.

\emph{Step~2: Combine contraction with
score--Jacobian relation.}
Proposition~\ref{prop:jac-contract} with
$u=aV(\alpha)$ and $v=bV(\beta)$
(for $a,b>0$ to be chosen) gives
\[
  \bigl\|J_{\fp}(aV(\alpha),bV(\beta))\bigr\|^2
  \le a^2\,\Phi_n(p)+b^2\,\Phi_n(q).
\]
By Proposition~\ref{prop:score-jac},
the left-hand side equals
$(a+b)^2\,\Phi_n(r)$.
Hence
\begin{equation}\label{eq:ab-ineq}
  (a+b)^2\,\Phi_n(r)
  \;\le\;a^2\,\Phi_n(p)+b^2\,\Phi_n(q).
\end{equation}

\emph{Step~3: Blachman's choice.}
Set $a=1/\Phi_n(p)$ and $b=1/\Phi_n(q)$.
Then~\eqref{eq:ab-ineq} becomes
\[
  \Bigl(\frac{1}{\Phi_n(p)}+\frac{1}{\Phi_n(q)}\Bigr)^{\!2}
  \Phi_n(r)
  \;\le\;\frac{1}{\Phi_n(p)}+\frac{1}{\Phi_n(q)}.
\]
Dividing both sides by
$\bigl(1/\Phi_n(p)+1/\Phi_n(q)\bigr)\Phi_n(r)>0$
yields~\eqref{eq:stam-full}.
\end{proof}


% ═══════════════════════════════════════════════════════════
\section{Discussion}\label{sec:discussion}
% ═══════════════════════════════════════════════════════════

The finite free Stam inequality
(Theorem~\ref{thm:stam-full}) is now established
for all $n\ge 2$.
The proof combines the commutation of $\fp$ with
differential operators, variance additivity, and the
convexity theory of hyperbolic polynomials.
The structural identities developed in
Section~\ref{sec:identities}
(Fisher--repulsion, Laplacian, Bezoutian)
and the Hermite heat equation
(Theorem~\ref{thm:hermite-heat}) provide
complementary perspectives.
The independent proofs for $n\le 3$
(Theorem~\ref{thm:n3-stam})
and for Gaussian inputs
(Theorem~\ref{thm:stam-gauss})
use different techniques
(sum-of-squares and dissipation, respectively)
and yield additional information
such as exact defect formulas.

\begin{remark}[Equality characterisation]
For $n=3$, Theorem~\ref{thm:n3-stam} shows
that equality in the Stam inequality holds
(for inputs with $u_p,u_q>0$)
if and only if $\ell_3(p)=\ell_3(q)=0$,
i.e., both inputs are finite Gaussians.
The characterisation of equality for
general~$n$ remains an interesting open question.
\end{remark}


% ═══════════════════════════════════════════════════════════
\section*{References}
% ═══════════════════════════════════════════════════════════

\begin{thebibliography}{99}

\bibitem{BGLS}
H.~H.~Bauschke, O.~G\"uler, A.~S.~Lewis,
and H.~S.~Sendov,
\emph{Hyperbolic polynomials and convex analysis},
Canad.\ J.\ Math.\ \textbf{53} (2001), 470--488.

\bibitem{bb}
J.~Borcea and P.~Br\"and\'en,
\emph{The {L}ee--{Y}ang and {P}\'olya--{S}chur programs.
{I}.\ {L}inear operators preserving stability},
Invent.\ Math.\ \textbf{177} (2009), 541--569.

\bibitem{blachman}
N.~M.~Blachman,
\emph{The convolution inequality for entropy powers},
IEEE Trans.\ Inform.\ Theory \textbf{IT-11} (1965),
267--271.

\bibitem{dembo}
A.~Dembo, T.~M.~Cover, and J.~A.~Thomas,
\emph{Information-theoretic inequalities},
IEEE Trans.\ Inform.\ Theory \textbf{37} (1991),
1501--1518.

\bibitem{dyson}
F.~J.~Dyson,
\emph{A {B}rownian-motion model for the eigenvalues
of a random matrix},
J.~Math.\ Phys.\ \textbf{3} (1962), 1191--1198.

\bibitem{fiedler}
M.~Fiedler,
\emph{Hankel and {L}oewner matrices},
Linear Algebra Appl.\ \textbf{58} (1984), 75--95.

\bibitem{GVS}
J.~Garza~Vargas, N.~Srivastava, and Z.~Stier,
\emph{The finite free {S}tam inequality},
preprint, 2026.

\bibitem{MSS1}
A.~Marcus, D.~A.~Spielman, and N.~Srivastava,
\emph{Interlacing families {I}: Bipartite {R}amanujan
graphs of all degrees},
Ann.\ of Math.\ (2) \textbf{182} (2015), 307--325.

\bibitem{MSS2}
A.~Marcus, D.~A.~Spielman, and N.~Srivastava,
\emph{Interlacing families {II}: Mixed characteristic
polynomials and the {K}adison--{S}inger problem},
Ann.\ of Math.\ (2) \textbf{182} (2015), 327--350.

\bibitem{nuij}
W.~Nuij,
\emph{A note on hyperbolic polynomials},
Math.\ Scand.\ \textbf{23} (1968), 69--72.

\bibitem{stam}
A.~J.~Stam,
\emph{Some inequalities satisfied by the quantities of
information of {F}isher and {S}hannon},
Inform.\ Control \textbf{2} (1959), 101--112.

\bibitem{voiculescu}
D.~V.~Voiculescu, K.~J.~Dykema, and A.~Nica,
\emph{Free Random Variables},
CRM Monograph Series, vol.~1, Amer.\ Math.\ Soc.,
Providence, RI, 1992.

\end{thebibliography}

\end{document}
